%% Figure environments for the quantmsdiann MCP manuscript.

%% Fig. 1 — pipeline architecture (subway, vector PDF from rsvg-convert)
%% on the left, runtime scaling panels stacked on the right.
\begin{figure}[htbp]
  \centering
  \begin{minipage}{\textwidth}
    \centering
    \begin{minipage}[t]{0.62\linewidth}
      \centering
      \textbf{(a)}\par\vspace{0.2em}
      \mcpincludegraphics[width=\linewidth,keepaspectratio]{manuscript/workflow_panel}
    \end{minipage}\hfill
    \begin{minipage}[t]{0.36\linewidth}
      \centering
      \textbf{(b)}\par\vspace{0.2em}
      \mcpincludegraphics[width=\linewidth,keepaspectratio]{performance/queue_size_sweep}\par\vspace{1.0em}
      \textbf{(c)}\par\vspace{0.2em}
      \mcpincludegraphics[width=\linewidth,keepaspectratio]{performance/parallelism_vs_wallclock}
    \end{minipage}
  \end{minipage}
  \caption[Pipeline and scaling]{%
    \textbf{\quantmsdiann{} pipeline architecture and runtime scaling.}
    (\textbf{a})~Subway diagram of the \quantmsdiann{} workflow: mandatory
    modules as filled circles, optional/skippable modules as open circles.
    Per-file parallel stages (red; file preparation, preliminary
    calibration, individual search) run as one task per raw file, whereas
    collective stages (green; in-silico library generation, empirical
    library assembly, final quantification, MSstats conversion, and
    \texttt{pmultiqc} reporting) run once over the full cohort.
    (\textbf{b})~Workflow wall-clock time versus the number of cluster nodes
    (Nextflow \texttt{executor.queueSize}) on the PXD071075 single-cell
    cohort.
    (\textbf{c})~Wall-clock time to finish each reanalysis, one bar per
    dataset, ordered by cohort size (number of MS data files) and coloured
    by instrument family. In both~(\textbf{b}) and~(\textbf{c}) the reported
    time \emph{excludes} the in-silico spectral-library prediction step
    (\texttt{INSILICO\_LIBRARY\_GENERATION}), a one-time, cohort-independent
    cost that can be predicted once and reused, so the bars reflect the
    per-file quantification work. Time-to-finish stays within a few hours
    across two orders of magnitude in file count because per-file stages are
    parallelised across the cluster; the deeper variable-modification search
    of the phosphoproteomics cohorts (e.g.\ PXD049692) is the main residual
    driver of wall-clock once library prediction is set aside.
  }
  \label{fig:architecture}
\end{figure}

%% Fig. 2 — ProteoBench: precursors-vs-version (compact summary) +
%% ID-vs-accuracy scatter (community context). The per-version per-module
%% protein-group / precursor breakdown is in Additional file 1 (Fig. S3).
\begin{figure}[htbp]
  \centering
  \mcpincludegraphics[width=\linewidth,keepaspectratio]{manuscript/fig2_proteobench_equivalence}
  \caption[quantmsdiann reproduces single-machine ProteoBench submissions]{%
    \textbf{Distributed \quantmsdiann{} reproduces single-machine results on
    ProteoBench.} For the two DIA modules with public predicted-from-FASTA
    community submissions (Orbitrap Astral, Module~7; timsTOF diaPASEF,
    PXD062685), \quantmsdiann{} (\texttt{DIA-NN} 1.8.1 and 2.5.1-enterprise) is
    compared against the cloud of community submissions, which are
    single-machine sequential runs. \quantmsdiann{} preserves
    match-between-runs by design: a preliminary pass builds a combined
    empirical library across all runs, against which each run is re-quantified.
    (\textbf{a})~Accuracy per concentration --- measured log$_2$ fold-change for
    each HYE species against the ProteoBench-expected ratio (dashed); community
    submissions as grey points.
    (\textbf{b})~Depth versus accuracy --- precursors quantified against
    median-$|\epsilon|$; \quantmsdiann{} (coloured) sits within the
    single-machine community cloud (grey), reaching greater depth at comparable
    accuracy from 1.8.1 to 2.5.1-enterprise.
    Datasets without predicted-library community comparators (single-cell
    Module~9/PXD049412; ZenoTOF Module~10/PXD070049) and the per-version
    protein-group counts are in Additional file~1.
  }
  \label{fig:proteobench}
\end{figure}

%% Fig. 3 — Independent cell-line reanalyses (depth vs original)
\begin{figure}[htbp]
  \centering
  \begin{minipage}{\textwidth}
    \centering
    \begin{minipage}[t]{0.32\linewidth}
      \centering
      {\scriptsize\textbf{(a)~PXD003539 (NCI-60, 2019)}}\par\vspace{0.2em}
      \mcpincludegraphics[width=\linewidth,keepaspectratio]{PXD003539/main_comparison}
    \end{minipage}\hfill
    \begin{minipage}[t]{0.32\linewidth}
      \centering
      {\scriptsize\textbf{(b)~MSV000093870 (plexDIA, 2024)}}\par\vspace{0.2em}
      \mcpincludegraphics[width=\linewidth,keepaspectratio]{plexDIA/MSV000093870/main_galatidou_total_pg}
    \end{minipage}\hfill
    \begin{minipage}[t]{0.32\linewidth}
      \centering
      {\scriptsize\textbf{(c)~PXD030304 (ProCan, 2022)}}\par\vspace{0.2em}
      \mcpincludegraphics[width=\linewidth,keepaspectratio]{PXD030304/main_comparison}
    \end{minipage}

    \vspace{1.0em}
    \begin{minipage}[t]{0.32\linewidth}
      \centering
      {\scriptsize\textbf{(d)~PXD064049 (MYCN DVP spatial, 2025)}}\par\vspace{0.2em}
      \mcpincludegraphics[width=\linewidth,keepaspectratio]{PXD064049/main_comparison}
    \end{minipage}\hfill
    \begin{minipage}[t]{0.32\linewidth}
      \centering
      {\scriptsize\textbf{(e)~PXD049692 (NK phospho, 2024)}}\par\vspace{0.2em}
      \mcpincludegraphics[width=\linewidth,keepaspectratio]{PXD049692/main_comparison}
    \end{minipage}\hfill
    \begin{minipage}[t]{0.32\linewidth}~\end{minipage}
  \end{minipage}
  \caption[Cell-line, spatial, and phospho reanalyses]{%
    \textbf{Independent reanalysis of public DIA datasets, spanning bulk cell
    lines, single cells, spatial proteomics, and phosphoproteomics, recovers
    comparable or greater identification depth than the original deposits.}
    Each panel compares the original published counts (grey) with the
    \quantmsdiann{} reanalysis at 1\% FDR (blue) from the same raw data.
    (\textbf{a})~NCI-60 bulk cell lines (PXD003539, Guo 2019; original
    OpenSWATH analysis).
    (\textbf{b})~Single-cell oocyte plexDIA (MSV000093870, Galatidou 2024).
    (\textbf{c})~ProCan-DepMapSanger pan-cancer cell-line panel (PXD030304).
    (\textbf{d})~MYCN-amplified neuroblastoma Deep Visual Proteomics
    (spatial, laser-microdissected regions; PXD064049, diaPASEF), with
    precursors and protein groups compared against the original
    DIA-NN~1.8.1 analysis.
    (\textbf{e})~NK-cell Fe-NTA phosphoproteomics (PXD049692, diaPASEF):
    stripped-sequence phosphopeptides (backbones) compared against the
    original Spectronaut directDIA analysis on the same runs.
    Per-cohort detail for panels~(a)--(c) in Additional file~1: Figs.~S4--S12.
  }
  \label{fig:cell-line-reanalysis}
\end{figure}

%% Fig. 4 — Pan-cohort of five DIA reanalyses
\begin{figure}[htbp]
  \centering
  \mcpincludegraphics[width=\textwidth,height=\mcpmaxfigheight,keepaspectratio]{manuscript/fig4_cellline_atlas}
  \caption[Pan-cohort of DIA reanalyses]{%
    \textbf{Pan-cohort of five public DIA cell-line reanalyses by
    \quantmsdiann{}.}
    Combines PXD003539 (NCI-60, Guo 2019), PXD030304
    (ProCan-DepMapSanger, 949 lines), PXD004701 (Sun 2023, 76 breast
    cancer lines), PXD041421 (Wang 2023), and PXD017199 (Tognetti
    2021, 67 breast cancer lines), reanalysed under a single
    \quantmsdiann{} invocation.
    (\textbf{a})~Per-cohort headline protein-group counts: the originally
    published count (grey) and the \quantmsdiann{} reanalysis (blue);
    PXD017199 and PXD041421 have no published DIA headline, so only the
    \quantmsdiann{} bar is shown.
    (\textbf{b})~Cell lines per pan-cancer tissue, stacked by contributing
    cohort.
    (\textbf{c})~Unique UniProt proteins detected per tissue across the
    combined panel.
  }
  \label{fig:pancohort}
\end{figure}
